\chapter{État de l'art}
\label{chap:etat_art}

\section{Biologie du Diabète de Type 1}

\subsection{Physiopathologie générale}
Le DT1 est caractérisé par une destruction des cellules bêta pancréatiques productrices d'insuline. [À COMPLÉTER : mécanismes auto-immuns, génétique HLA, facteurs environnementaux]

\subsection{Spécificités africaines}
Formes idiopathiques (20-60\%), absence d'auto-anticorps, facteurs génétiques distincts \cite{Katte2023,Sobngwi2002}.

\section{Intelligence Artificielle en médecine}

\subsection{Machine Learning supervisé}
Classification, régression, métriques (recall, precision, AUC-ROC) \cite{Bishop2006,Hastie2009}.

\subsection{Algorithmes de classification}
Random Forest \cite{Breiman2001}, XGBoost \cite{Chen2016}, LightGBM \cite{Ke2017}, SVM.

\subsection{Interprétabilité}
SHAP \cite{Lundberg2017}, LIME \cite{Ribeiro2016}. Importance pour confiance clinique.

\section{Applications ML pour le DT1}

\subsection{Modèles existants}
Revue des systèmes de prédiction, scores de risque génétique \cite{Contreras2018}.

\subsection{Limitations en Afrique}
Dataset shift \cite{Quinonero-Candela2009}, manque de données locales, performances dégradées \cite{Oram2025}.

\section{Positionnement de notre approche}
Notre travail se distingue par : (1) focus populations camerounaises, (2) biomarqueurs validés, (3) interprétabilité prioritaire.
